\documentclass[11pt,letterpaper]{article}

\usepackage[T1]{fontenc}
\usepackage[utf8]{inputenc}
\usepackage[french]{babel}
\usepackage[margin=1.7cm]{geometry}
\usepackage{array}
\usepackage{booktabs}
\usepackage{enumitem}
\usepackage{fancyhdr}
\usepackage{lastpage}
\usepackage{lmodern}
\usepackage{longtable}
\usepackage{ragged2e}
\usepackage{tabularx}
\usepackage{tikz}
\usepackage[table]{xcolor}
\usepackage{hyperref}

\renewcommand{\familydefault}{\sfdefault}
\setlength{\parindent}{0pt}
\setlength{\parskip}{0.45em}
\setlength{\LTpre}{0.35em}
\setlength{\LTpost}{0.35em}
\setlist[itemize]{label=\textbullet,leftmargin=1.2em,itemsep=0.12em,topsep=0.15em}
\setlist[enumerate]{leftmargin=1.3em,itemsep=0.12em,topsep=0.15em}

\definecolor{SequoInk}{HTML}{1F2933}
\definecolor{SequoMuted}{HTML}{59636F}
\definecolor{SequoGreen}{HTML}{1F7A5B}
\definecolor{SequoBlue}{HTML}{164E63}
\definecolor{SequoPale}{HTML}{F3F7F5}
\definecolor{SequoLine}{HTML}{D9E2DE}
\definecolor{SequoHeader}{HTML}{E9F1EE}
\definecolor{GaboOlive}{HTML}{6D5E0F}
\definecolor{GaboCream}{HTML}{FFF9EE}
\arrayrulecolor{SequoLine}
\newcolumntype{Y}{>{\RaggedRight\arraybackslash}X}
\newcolumntype{L}[1]{>{\RaggedRight\arraybackslash}p{#1}}

\hypersetup{
  colorlinks=true,
  urlcolor=SequoBlue,
  linkcolor=SequoBlue
}

\pagestyle{fancy}
\fancyhf{}
\renewcommand{\headrulewidth}{0pt}
\fancyfoot[L]{\small\color{SequoMuted}Compte rendu mensuel - Sequo - Septembre 2026}
\fancyfoot[R]{\small\color{SequoMuted}Page \thepage\ sur \pageref{LastPage}}

\newcommand{\gabologo}{%
  \begin{tikzpicture}[x=0.85pt,y=0.85pt]
    \fill[GaboCream, rounded corners=32pt] (0,0) rectangle (150,150);
    \draw[color=GaboOlive,line width=13pt,line join=round,line cap=round]
      (106,108) -- (64,108) -- (37,75) -- (64,42) -- (106,42);
    \draw[color=GaboOlive,line width=13pt,line cap=round] (88,75) -- (124,75);
    \fill[GaboOlive] (126,75) circle (6.5pt);
  \end{tikzpicture}%
}

\newcommand{\sectiontitle}[1]{%
  \vspace{0.55em}
  {\Large\bfseries\color{SequoInk}#1}\par
  \vspace{0.1em}
}

\newcommand{\subsectiontitle}[1]{%
  \vspace{0.35em}
  {\large\bfseries\color{SequoInk}#1}\par
}

\newcommand{\simpleexplanation}[1]{%
  {\small\textbf{\color{SequoBlue}Explication simple.} \color{SequoMuted}#1\par}
}

\newcommand{\metric}[2]{%
  \begin{minipage}[t]{0.30\linewidth}
    {\LARGE\bfseries\color{SequoGreen}#1}\par
    {\small\color{SequoMuted}#2}
  \end{minipage}
}

\newcommand{\statusdone}{\textbf{\color{SequoGreen}Réalisé}}
\newcommand{\statuspartial}{\textbf{\color{SequoBlue}En cours}}
\newcommand{\statuspending}{\textbf{\color{SequoMuted}À finaliser}}

\begin{document}
\pagenumbering{roman}
\color{SequoInk}
\thispagestyle{empty}

\begin{titlepage}
  \thispagestyle{empty}
  \begin{tikzpicture}[remember picture,overlay]
    \fill[SequoPale] (current page.north west) rectangle ([yshift=-4.8cm]current page.north east);
    \fill[GaboOlive] (current page.north west) rectangle ([xshift=0.38cm]current page.south west);
    \fill[SequoHeader] ([yshift=4.1cm]current page.south west) rectangle (current page.south east);
  \end{tikzpicture}

  \vspace*{0.7cm}
  \begin{minipage}[t]{0.52\linewidth}
    \resizebox{1.65cm}{!}{\gabologo}\hspace{0.35cm}%
    \raisebox{0.48cm}{\Large\bfseries\color{GaboOlive}GABO}
  \end{minipage}
  \hfill
  \begin{minipage}[t]{0.38\linewidth}
    \raggedleft
    {\small\bfseries RAPPORT MENSUEL}\par
    {\small Écosystème Sequo}\par
    {\small\color{SequoMuted}30 septembre 2026}
  \end{minipage}

  \vspace{2.0cm}
  {\small\bfseries\color{SequoMuted}Période de suivi}\par
  {\LARGE\bfseries\color{GaboOlive}Septembre 2026}\par
  \vspace{0.45cm}
  {\fontsize{31}{34}\selectfont\bfseries Compte rendu mensuel\\du travail accompli}\par
  \vspace{0.35cm}
  {\Large Projet Sequo}\par
  \vspace{0.15cm}
  {\large\color{SequoMuted}API, application client, SequoHub, sites web public, admin et développeurs}\par

  \vfill
  {\renewcommand{\arraystretch}{1.22}
  \begin{tabularx}{0.88\linewidth}{@{}L{4.0cm}Y@{}}
    \toprule
    \textbf{Période couverte} & Du 1er au 30 septembre 2026 \\
    \textbf{Destinataire} & Gamal Tchadenu / SEQUO \\
    \textbf{Préparé par} & Muhirwa Gabo Oreste \\
    \textbf{Objet} & Bilan d'avancement du mois de septembre 2026 \\
    \textbf{Base de vérification} & Historique Git des dépôts Sequo consultés localement \\
    \textbf{Format} & Rapport technique avec explications simples \\
    \bottomrule
  \end{tabularx}}

  \vspace{0.55cm}
  {\small
  \textbf{Notice de lecture.} Ce rapport conserve les informations techniques nécessaires au suivi du projet. Les passages intitulés \textit{Explication simple} reformulent certains points pour un lecteur non développeur. Ces explications sont volontairement moins complètes et moins précises que les détails techniques : elles servent à faciliter la compréhension, pas à remplacer les spécifications ou le code.
  }
\end{titlepage}

\pagenumbering{arabic}
\pagestyle{fancy}

\sectiontitle{Périmètre du rapport}
Ce compte rendu couvre uniquement les travaux enregistrés dans les dépôts Git pendant le mois de septembre 2026. Il ne reprend pas les travaux d'avant septembre déjà présentés dans le rapport précédent.

\rowcolors{2}{SequoPale}{white}
{\small
\begin{tabularx}{\linewidth}{@{}L{3.1cm}L{4.5cm}Y@{}}
\toprule
\rowcolor{SequoHeader}
\textbf{Projet} & \textbf{Type} & \textbf{Lecture du mois} \\
\midrule
Sequo API & Backend / API & Plus gros avancement technique du mois, avec sécurité, commandes, livraisons, notifications, realtime, documentation et mise en ligne sur \href{https://api.sequoservice.com}{api.sequoservice.com}. \\
Sequo & Application client & Avancement important sur l'interface mobile, le panier, le catalogue, les notifications et les parcours de commande. Les données affichées restent fictives et ne viennent pas encore de l'API. \\
SequoHub & Application point relais & Avancement important sur l'application de comptoir : casiers, scan, remise, historique, préférences et premières bases d'authentification. L'intégration réelle des données API reste à faire. \\
sequo-services & Site public & Mise en place et déploiement du site \href{https://www.sequoservice.com}{www.sequoservice.com}. Le contenu est encore de type placeholder / hello world et doit être remplacé. \\
sequo-admin & Site admin & Mise en place et déploiement du site \href{https://admin.sequoservice.com}{admin.sequoservice.com}. Le contenu et les données sont encore fictifs ou de démonstration. \\
sequo-developers & Site développeurs & Mise en place et déploiement du site \href{https://developers.sequoservice.com}{developers.sequoservice.com}. La structure est prête, mais le contenu est encore provisoire et incomplet. \\
\bottomrule
\end{tabularx}
}

\simpleexplanation{En simple, septembre a servi à faire avancer le serveur principal, deux applications mobiles et les trois sites web. Les applications utilisent encore des données fictives, pas les données réelles de l'API. Les sites sont déjà en ligne, mais leur contenu actuel est encore provisoire, proche d'un placeholder / hello world : l'objectif du mois était surtout d'établir les technologies, le déploiement et les bases de chaque portail.}

\sectiontitle{Résumé exécutif}
Le mois de septembre 2026 a été un mois de construction important pour l'écosystème Sequo. Le travail le plus conséquent a été réalisé sur Sequo API, SequoHub et Sequo. En parallèle, les trois sites web ont été créés, structurés et mis en ligne afin de disposer des domaines publics, du socle technique et des premières interfaces.

Le backend Sequo API a fortement avancé : sécurité, premières bases d'authentification Google, sessions, règles de production, catalogues, commandes, paiements, retours, relay, casiers, notifications, WebSocket, provider delivery, documentation API et procédures de déploiement. Le système d'authentification est toutefois commencé mais pas terminé. Les applications Sequo et SequoHub ne sont pas terminées : elles utilisent encore des données fictives, sans intégration réelle des données venant de l'API, mais elles ont gagné des écrans concrets, des composants UI et plusieurs workflows métiers.

\simpleexplanation{En langage simple, le mois a surtout permis de transformer Sequo en un ensemble plus réel : le serveur est accessible en ligne, les sites existent sur leurs domaines, et les applications commencent à avoir des écrans proches du produit final. Mais les applications ne consomment pas encore les vraies données de l'API, l'authentification n'est pas finalisée, et les sites doivent encore recevoir leur vrai contenu.}

\vspace{0.45em}
\noindent
\metric{1\,096}{commits enregistrés en septembre}
\hfill
\metric{6}{dépôts suivis dans ce rapport}
\hfill
\metric{4}{services déjà accessibles en ligne}
\par\vspace{0.65em}
\noindent
\metric{692}{commits Sequo API}
\hfill
\metric{193}{commits SequoHub}
\hfill
\metric{101}{commits Sequo}

\sectiontitle{Repères non techniques}
Ce tableau traduit les termes informatiques les plus utiles pour lire le rapport.

\rowcolors{2}{SequoPale}{white}
{\small
\begin{tabularx}{\linewidth}{@{}L{3.4cm}Y Y@{}}
\toprule
\rowcolor{SequoHeader}
\textbf{Terme} & \textbf{Explication simple} & \textbf{Pourquoi c'est important pour Sequo} \\
\midrule
API / backend & Le serveur qui garde les règles, les données et les décisions importantes. & C'est le moteur commun utilisé par les applications mobiles et les sites internes. \\
Déploiement & Le fait de rendre un service accessible sur internet avec son domaine. & Cela permet de tester hors ordinateur local et de préparer les vraies intégrations. \\
Vue / Vite & Technologies utilisées pour construire rapidement les sites web. & Elles permettent de créer les portails public, admin et développeurs avec une base moderne. \\
Kotlin Multiplatform & Technologie utilisée pour partager du code entre Android et iOS. & Elle aide à construire les applications mobiles plus vite en évitant de réécrire deux fois certaines parties. \\
Google Sign-In & Connexion avec un compte Google. & Utile pour simplifier l'accès aux applications et réduire les erreurs de mot de passe. \\
FCM / notifications push & Service qui permet d'envoyer des alertes sur mobile. & Nécessaire pour prévenir clients, relay, livreurs ou opérateurs au bon moment. \\
WebSocket & Canal de communication en direct entre l'application et le serveur. & Utile pour suivre des événements en temps réel : mission, commande, négociation, notification. \\
Outbox & File d'attente interne pour envoyer des événements de façon fiable. & Évite d'envoyer des notifications ou actions si l'opération principale n'a pas vraiment réussi. \\
Provider delivery & Préparation des envois vers des services externes, par exemple push ou SMS. & Permet de gérer les essais, erreurs, reprises et audits d'envoi. \\
\bottomrule
\end{tabularx}
}

\newpage
\sectiontitle{Travail accompli par projet}

\subsectiontitle{1. Sequo API -- backend principal}
Le backend est le chantier le plus avancé du mois. Il est maintenant associé au domaine public \href{https://api.sequoservice.com}{api.sequoservice.com}, avec documentation de déploiement, configuration production, Docker, healthcheck et notes VPS.

\begin{itemize}
  \item Avancement du chantier d'authentification : Google OAuth, validation des tokens Google, gestion des erreurs, liaison de comptes existants, sessions persistantes, refresh sessions et tests associés. Le système est commencé, mais pas encore terminé pour une utilisation complète en production.
  \item Mise en place et durcissement de la configuration de production : variables d'environnement, secrets, CORS explicite, garde-fous contre les valeurs de test, Docker Compose, documentation VPS et endpoint public.
  \item Avancement important des domaines métier : commandes, panier, paiements, webhooks, commissions, settlements, retours, relay parcels, pickup codes, casiers, heures d'ouverture Hub et règles opérationnelles.
  \item Ajout des bases de notifications : inbox, préférences, routing, outbox, worker, delivery planning, provider delivery, WebSocket delivery et suivi de présence.
  \item Ajout de WebSocket/STOMP avec authentification JWT, autorisation des souscriptions, présence des clients et publication d'événements realtime pour notifications, missions rider et bargaining.
  \item Création ou mise à jour de nombreux documents techniques : guide mobile API, diagrammes PlantUML, Postman endpoints, architecture auth, paiements, notifications, déploiement et notes VPS.
\end{itemize}

\simpleexplanation{En simple, le serveur est beaucoup plus proche d'un socle exploitable, mais la connexion utilisateur n'est pas encore un chantier terminé. Certaines parties existent et sont testées, mais il reste du travail avant de considérer l'authentification comme finalisée pour la production.}

\subsectiontitle{2. Sequo -- application client}
L'application client a reçu une progression visible sur l'expérience utilisateur. Elle prépare la future connexion au backend, mais les écrans utilisent encore des données fictives et ne consomment pas encore les données réelles de l'API.

\begin{itemize}
  \item Construction des écrans marketplace : accueil, catégories, boutiques, fiches produits, galerie photo, recommandations et produits alternatifs.
  \item Amélioration du panier, des commandes et des états de commande, avec une interface plus lisible pour les parcours actuels et passés.
  \item Préparation du parcours d'authentification, Google Sign-In, stockage sécurisé de session, configuration réseau et client HTTP Ktor. L'intégration complète avec les endpoints réels reste à faire.
  \item Ajout d'une section notifications avec inbox, filtres, archivage, restauration, actions rapides, états vides et visibilité contrôlée des codes de retrait.
  \item Amélioration du design system : thèmes, icônes, formes, composants d'erreur, top app bar, navigation drawer, dark mode et assets.
  \item Nettoyage de certaines anciennes logiques locales pour préparer une application davantage pilotée par l'API plus tard.
\end{itemize}

\simpleexplanation{En simple, l'application client commence à ressembler à un vrai produit mobile côté écran et navigation. Mais elle fonctionne encore avec des données fictives : les données réelles de l'API ne sont pas encore intégrées.}

\subsectiontitle{3. SequoHub -- application point relais}
SequoHub a fortement avancé en septembre. Le dépôt contient maintenant une base concrète pour l'application utilisée au comptoir des points relais.

\begin{itemize}
  \item Mise en place du dashboard Hub : suivi des casiers, états d'attention, recherche, tabs, historique et indicateurs visuels.
  \item Ajout des écrans de réception, scan, remise client, retours, paramètres, préférences, notifications, activité et gestion des heures d'ouverture.
  \item Préparation de l'authentification : Google Sign-In Android/iOS, stockage de session, client HTTP, health check API, gestion des erreurs et états de chargement. Le parcours complet reste à finaliser.
  \item Mise en place de la localisation et de la langue, avec refactoring progressif des textes affichés.
  \item Ajout d'assets de marque : logo, splash screen, icônes Android, onboarding et documentation de branding.
  \item Alignement avec l'API mobile : mapping des endpoints SequoHub, TODO mobile API, guides opérationnels et documentation relay. Les écrans restent majoritairement alimentés par des données fictives.
\end{itemize}

\simpleexplanation{En simple, SequoHub est l'application du point relais. Le mois de septembre a permis de créer la majorité de son squelette visible : casiers, scan, remise, historique, connexion et paramètres. Elle n'est pas encore prête terrain et n'utilise pas encore les vraies données API dans ses écrans principaux.}

\subsectiontitle{4. Sites web mis en ligne}
Les sites web ont été créés et déployés en septembre. Le contenu actuel est encore très provisoire, de type placeholder / hello world, et ne correspond pas à la version commerciale finale. Le but principal était d'établir les domaines, la technologie, les builds et les premières interfaces.

\rowcolors{2}{SequoPale}{white}
{\small
\begin{tabularx}{\linewidth}{@{}L{3.1cm}L{4.6cm}Y@{}}
\toprule
\rowcolor{SequoHeader}
\textbf{Site} & \textbf{URL} & \textbf{Travail de septembre} \\
\midrule
Site public & \href{https://www.sequoservice.com}{www.sequoservice.com} & Projet Vue 3 + Vite + TypeScript, structure de landing page, styles globaux, branding, documentation produit, roadmap et guide de développement. Le contenu en ligne reste placeholder / hello world et doit être remplacé. \\
Site admin & \href{https://admin.sequoservice.com}{admin.sequoservice.com} & Projet Vue 3 + Vite + TypeScript, shell administratif, navigation, vues opérationnelles avec données fictives, thème clair/sombre, documentation UX, sécurité et roadmap. Le site n'est pas encore connecté aux vraies données opérationnelles. \\
Site développeurs & \href{https://developers.sequoservice.com}{developers.sequoservice.com} & Projet Vue 3 + Vite + TypeScript, portail documentaire interactif, architecture de documentation API, guide d'écriture, tech stack, roadmap et base pour futures intégrations. Le contenu reste provisoire et incomplet. \\
\bottomrule
\end{tabularx}
}

\simpleexplanation{Ces sites sont en ligne pour poser les fondations : domaine, hébergement, technologie et première structure. Leur contenu actuel ne doit pas être considéré comme final : il est encore de type placeholder / hello world et devra être repris avant communication large.}

\newpage
\sectiontitle{État synthétique}
\rowcolors{2}{SequoPale}{white}
{\small
\begin{tabularx}{\linewidth}{@{}L{3.2cm}L{2.0cm}Y@{}}
\toprule
\rowcolor{SequoHeader}
\textbf{Chantier} & \textbf{État} & \textbf{Commentaire} \\
\midrule
API publique & \statusdone & Le domaine API est documenté et l'environnement de déploiement a été préparé. Des validations production restent nécessaires avant ouverture complète. \\
Sécurité et auth & \statuspartial & Le système d'authentification est commencé mais pas terminé. Google Sign-In, validation tokens, sessions et guardrails ont avancé, mais il reste les parcours complets, audits, contrôles par rôle et validations production. \\
Commandes / paiements & \statuspartial & Les modèles, webhooks, checkouts, commissions et settlements ont fortement progressé. Les prestataires réels et la réconciliation production restent à finaliser. \\
Relay / Hub & \statuspartial & API Hub et application SequoHub ont beaucoup avancé. Les écrans utilisent encore surtout des données fictives et l'intégration réelle avec l'API doit continuer. \\
Notifications / realtime & \statuspartial & Inbox, outbox, routing, provider delivery et WebSocket ont avancé. L'envoi réel à grande échelle et les workflows complets restent à stabiliser. \\
Application client & \statuspartial & Les écrans principaux et le socle auth/réseau existent. Les données affichées restent fictives et la connexion complète aux endpoints réels reste le prochain grand travail. \\
Sites web & \statuspartial & Les trois sites sont en ligne et technologiquement établis. Le contenu est encore placeholder / hello world et doit être repris. \\
Documentation & \statusdone & Guides API, diagrammes, notes de déploiement et documents d'architecture ont été fortement enrichis. \\
\bottomrule
\end{tabularx}
}

\sectiontitle{Travail non finalisé}
Les éléments suivants restent ouverts après septembre. Ils sont importants pour la suite, mais ils ne remettent pas en cause le travail accompli pendant le mois.

\rowcolors{2}{SequoPale}{white}
{\footnotesize
\begin{longtable}{@{}L{3.4cm}L{6.4cm}L{6.4cm}@{}}
\toprule
\rowcolor{SequoHeader}
\textbf{Sujet} & \textbf{Ce qui reste à faire} & \textbf{Explication simple / impact} \\
\midrule
\endfirsthead
\toprule
\rowcolor{SequoHeader}
\textbf{Sujet} & \textbf{Ce qui reste à faire} & \textbf{Explication simple / impact} \\
\midrule
\endhead
Contenu des sites & Reprendre textes, images, preuves, arguments commerciaux, pages légales et cohérence des domaines. & Les sites sont en ligne, mais leur contenu n'est pas encore assez bon pour représenter officiellement Sequo auprès du public. \\
Applications mobiles & Remplacer les données fictives par les données réelles venant de l'API, connecter les écrans aux endpoints réels, gérer les erreurs réseau, les états vides, la persistance et les tests mobiles. & Les applications existent et avancent visuellement, mais elles ne consomment pas encore les vraies données de l'API. \\
Paiements réels & Finaliser les intégrations Yas Togo / Moov Africa, callbacks signés, idempotence, tests et réconciliation. & Les règles sont préparées, mais les vrais paiements demandent une intégration contrôlée avec les prestataires. \\
Notifications production & Finaliser provider push/SMS, budgets, retries, monitoring et scénarios par rôle. & Les notifications doivent arriver à la bonne personne, au bon moment, sans doublons ni pertes. \\
Temps réel mobile & Continuer la validation WebSocket côté applications, présence, droits par topic et publication des événements métier. & Le direct est utile pour les livreurs, hubs et commandes, mais il doit être sécurisé et fiable. \\
Admin réel & Remplacer les données fictives par l'API, finaliser l'authentification, les rôles, les tables, les filtres et les actions admin contrôlées. & Le site admin est une base visuelle, pas encore l'outil opérationnel complet. \\
Qualité production & Compléter smoke tests, monitoring, secrets, sauvegardes, logs, alertes et procédures d'incident. & Avant ouverture large, il faut pouvoir détecter, comprendre et corriger rapidement les incidents. \\
\bottomrule
\end{longtable}
}

\sectiontitle{Prochaines actions prévues}
Pour le prochain cycle, le travail prévu est de se concentrer sur les chantiers réellement prioritaires : connecter les applications aux endpoints existants, finaliser le parcours d'authentification, et continuer l'avancement de SequoHub, Sequo et Sequo API.

\begin{enumerate}
  \item Remplacer les données fictives des applications Sequo et SequoHub par des données venant des endpoints backend prioritaires.
  \item Finaliser le système d'authentification et les parcours de session, avec gestion propre des erreurs.
  \item Avancer l'application SequoHub sur les parcours point relais : scan, réception, remise, casiers, historique et états d'erreur.
  \item Avancer l'application Sequo sur les parcours client : catalogue, panier, commandes, notifications et navigation.
  \item Continuer l'avancement de Sequo API pour soutenir ces intégrations mobiles et corriger les endpoints manquants ou incomplets.
  \item Garder les autres chantiers, comme les paiements mobiles et l'amélioration complète des sites web, hors priorité principale pour ce cycle.
\end{enumerate}

\simpleexplanation{En simple, le prochain mois sera surtout consacré à relier les applications mobiles au serveur et à terminer la connexion utilisateur. Les paiements mobiles et la refonte complète des sites web ne sont pas annoncés comme priorités de ce cycle.}

\sectiontitle{Annexe de traçabilité}
Les chiffres ci-dessous proviennent de l'historique Git local des dépôts consultés pour la période du 1er au 30 septembre 2026.

\rowcolors{2}{SequoPale}{white}
{\small
\begin{tabularx}{\linewidth}{@{}L{3.3cm}L{2.0cm}Y@{}}
\toprule
\rowcolor{SequoHeader}
\textbf{Dépôt} & \textbf{Commits} & \textbf{Lecture synthétique} \\
\midrule
sequo-api & 692 & Backend, sécurité, paiements, notifications, realtime, Hub API, documentation, déploiement et diagrammes. \\
SequoHub & 193 & Application point relais, dashboard, casiers, scan, bases auth, API client, branding et localisation. Données applicatives encore fictives. \\
Sequo & 101 & Application client, marketplace, catalogue, panier, notifications, bases auth, design system et API client. Données applicatives encore fictives. \\
sequo-admin & 51 & Socle admin Vue/Vite, dashboard, navigation, thèmes, données de démonstration et documentation. \\
sequo-services & 30 & Site public Vue/Vite, structure marketing, styles, documentation produit et roadmap. \\
sequo-developers & 29 & Portail développeurs Vue/Vite, documentation API, architecture de contenu et guides. \\
\bottomrule
\end{tabularx}
}

\sectiontitle{Conclusion}
Septembre 2026 a été un mois de progression importante pour Sequo. Le backend a beaucoup avancé et dispose maintenant d'un domaine API public documenté. Le système d'authentification est commencé mais pas terminé. Les applications Sequo et SequoHub ont gagné des interfaces, une base réseau et plusieurs workflows visibles, mais elles utilisent encore des données fictives. Les trois sites web ont été créés et mis en ligne afin de poser les fondations publiques, administratives et développeurs, avec un contenu encore placeholder / hello world.

Le projet n'est pas encore terminé : l'authentification doit être finalisée, les applications doivent encore remplacer les données fictives par les données réelles de l'API, les contenus web doivent être repris, les paiements réels doivent être finalisés et les garanties de production doivent continuer à être renforcées. Le mois a néanmoins permis de passer d'un ensemble de dépôts en construction à un écosystème Sequo plus structuré, partiellement déployé et prêt à poursuivre vers des parcours utilisables.

\end{document}
